% Example of a two-column article with a compact preamble.
% All content is placeholder: lipsum text, mwe example images, and a
% dummy bibliography that is written out at compile time (filecontents).
%
% Build:  latexmk -pdf paper_2col_example.tex

\documentclass[10pt,twocolumn]{article}

\usepackage[margin=0.75in]{geometry}
\usepackage{lmodern}
\usepackage{microtype}
\usepackage{parskip}
\usepackage[compact]{titlesec}
\setcounter{secnumdepth}{0} % no section numbering

\usepackage{amsmath,amssymb}
\usepackage{graphicx}
\usepackage{tabularx,booktabs}

\usepackage[numbers,sort&compress]{natbib}
\bibliographystyle{IEEEtranN}

\usepackage[hidelinks]{hyperref}

% Float tuning: let large floats sit at the top of a text page instead of
% being pushed onto float-only pages
\renewcommand{\topfraction}{0.9}
\renewcommand{\dbltopfraction}{0.9}
\renewcommand{\textfraction}{0.07}
\renewcommand{\floatpagefraction}{0.85}
\renewcommand{\dblfloatpagefraction}{0.85}
\setlength{\emergencystretch}{3em} % prevent overfull lines

\usepackage{lipsum} % placeholder text

% Dummy bibliography, written next to this file when compiled.
\begin{filecontents*}[overwrite]{paper_2col_example.bib}
@article{Author2020,
  author  = {Author, First A. and Author, Second B.},
  title   = {A placeholder article title},
  journal = {Journal of Placeholder Studies},
  year    = {2020},
  volume  = {12},
  pages   = {1--10}
}
@article{Author2021,
  author  = {Author, Second B. and Author, Third C.},
  title   = {Another placeholder article title},
  journal = {Placeholder Letters},
  year    = {2021},
  volume  = {3},
  pages   = {100--110}
}
@inproceedings{Author2022,
  author    = {Author, Third C.},
  title     = {A placeholder conference paper},
  booktitle = {Proceedings of the Placeholder Conference},
  year      = {2022},
  pages     = {42--48}
}
@book{Author2019,
  author    = {Author, First A.},
  title     = {A Placeholder Book},
  publisher = {Placeholder Press},
  year      = {2019}
}
@misc{Author2023,
  author       = {Author, Second B.},
  title        = {A placeholder preprint},
  year         = {2023},
  howpublished = {arXiv:0000.00000}
}
\end{filecontents*}

\title{A Placeholder Title for an Example Two-Column Article}
\author{
First A.\ Author\thanks{Institution One} \and
Second B.\ Author\thanks{Institution Two} \and
Third C.\ Author\thanks{Institution Three}}
\date{\today}  % leave empty to suppress date, or put a date in

\begin{document}

\twocolumn[
\maketitle
\begin{quotation}\noindent\small
{\lipsum[1]} % braces needed: a bare ] would end the \twocolumn[...] argument
\end{quotation}
\vspace{1.5em}
]
% \thanks footnotes are lost inside \twocolumn[...], so re-issue them here
\renewcommand{\thefootnote}{\fnsymbol{footnote}}
\footnotetext[1]{Institution One}
\footnotetext[2]{Institution Two}
\footnotetext[3]{Institution Three}
\renewcommand{\thefootnote}{\arabic{footnote}}
\setcounter{footnote}{0}


\lipsum[2][1-4] Prior work has addressed this in several ways
\cite{Author2020,Author2021,Author2022}. \lipsum[2][5-8]

\lipsum[3]

\lipsum[4][1-3] We call this the \textbf{placeholder effect} (PE), following
earlier use of the term \cite{Author2019, Author2023}. \lipsum[4][4-9]

\lipsum[5]

\section{First section, with a single-column figure}\label{first-section}

\lipsum[6][1-5] This is illustrated in Figure \ref{fig:one}a. \lipsum[6][6-8]

\begin{figure}[htbp]
  \centering
  \includegraphics[width=.9\columnwidth]{example-image-a}
  \caption{
  Single-column figure placeholder.
  \textbf{A}. \lipsum[7][1-2]
  \textbf{B}. \lipsum[7][3-4]
  \textbf{C}. \lipsum[7][5-6]
  }
  \label{fig:one}
\end{figure}

\lipsum[8][1-3] In our tests at the time of writing\footnote{This is an
  example of a footnote in the body text.} we observed the effect
consistently (Figure \ref{fig:one}b). \lipsum[8][4-8]

\lipsum[9][1-4] If a task contains $n$ subtasks and the probability of
success for each subtask is $p$, the probability of success for the
overall task is $p^n$. \lipsum[9][5-7]

\lipsum[10]

\section{Second section, with a two-column table and figure}\label{second-section}

\lipsum[11]

\lipsum[12][1-4] Table \ref{tab:tiers} summarises the categories used
throughout. \lipsum[12][5-8]

\begin{table*}[htbp]
  \centering
  \begin{tabularx}{\linewidth}{l >{\raggedright\arraybackslash}p{2.5cm} >{\raggedright\arraybackslash}X}
    \toprule
    Tier & Level & Examples \\
    \midrule
    L1 & High &
      \lipsum*[13][1-2] \\
    L2 & Partial &
      \lipsum*[14][1-3] \\
    L3 & Low &
      \lipsum*[15][1-2] \\
    \bottomrule
  \end{tabularx}
  \caption{Placeholder table caption, adapted from \cite{Author2021}.
    \lipsum*[16][1-2]}
  \label{tab:tiers}
\end{table*}

\lipsum[17]

\lipsum[18]

\begin{figure*}[tb] % no p: keep tall wide figures off float-only pages
  \centering
  \includegraphics[width=\textwidth]{example-image-16x9}
  \caption{Two-column figure placeholder. \lipsum*[19][1-3]}
  \label{fig:two}
\end{figure*}

\section{Third section}\label{third-section}

\lipsum[20][1-6] See Figure \ref{fig:two}. \lipsum[20][7-9]

\lipsum[21]

\lipsum[22]

\begin{figure*}[tb] % no p: keep tall wide figures off float-only pages
  \centering
  \includegraphics[width=\textwidth]{example-image-golden}
  \caption{Another two-column figure placeholder
  \cite{Author2020,Author2022}. \lipsum*[23][1-3]}
  \label{fig:three}
\end{figure*}

\section{Fourth section}

\lipsum[24][1-5] Two clear exceptions are shown in Figure
\ref{fig:three} \cite{Author2023}. \lipsum[24][6-8]

\lipsum[25]

\section{Fifth section, with an itemised list}\label{fifth-section}

\lipsum[26]

\lipsum[27][1-3]

\begin{itemize}
\item
  \emph{\textbf{First item}}: \lipsum[28][1-4] \cite{Author2019}.
\item
  \emph{\textbf{Second item}}: \lipsum[29][1-4] \cite{Author2020,Author2021}.
\item
  \emph{\textbf{Third item}}: \lipsum[30][1-3]
\item
  \emph{\textbf{Fourth item}}: \lipsum[31][1-4] \cite{Author2022}.
\item
  \emph{\textbf{Fifth item}}: \lipsum[32][1-3]
\item
  \emph{\textbf{Sixth item}}: \lipsum[33][1-4] \cite{Author2023}.
\end{itemize}

\lipsum[34]

\lipsum[35][1-5] Prior experience suggests that \textasciitilde500 hours
of well-labelled data is a useful intermediate milestone
\cite{Author2021}. \lipsum[35][6-8]

\lipsum[36]

\section{Discussion}\label{discussion}

\lipsum[37]

\lipsum[38][1-4] We adopt the term deliberately to connect with earlier
work \cite{Author2019,Author2020,Author2023}. \lipsum[38][5-8]

\lipsum[39]

\section{Methods}\label{methods}

\textbf{First procedure.} \lipsum[40][1-5]

\textbf{Second procedure.} \lipsum[41][1-3] Each curve is a
two-parameter Weibull survival function $S(t) = 0.8^{(t/t_{80})^k}$,
plotted against $\log_{10}(t/t_{80})$, with $k \approx 0.68$ estimated
via $k = \ln(0.322)/\ln(0.19)$. Data were downloaded from
\href{https://example.org/}{example.org}. \lipsum[41][4-5]

\textbf{Third procedure.} \lipsum[42][1-5] \cite{Author2022}

\textbf{Fourth procedure.} \lipsum[43][1-5]

\bibliography{paper_2col_example}

\end{document}
